\section{ESP-IDF и JTAG: увидеть выполнение}\label{les:idf-jtag}
\lessonintro
  {использовать точки останова и воспроизводимую сборку}
  {уроки 14 и 28}
  {найденная причина целочисленного деления и исправленная функция}
  {ESP-IDF 5.4.x, OpenOCD Espressif, нативный USB}

\subsection{Когда печати недостаточно}
Отладчик позволяет остановиться перед строкой, посмотреть локальные переменные,
стек вызовов и пройти код по шагам. Breakpoint задаёт место остановки, watch показывает
значение, backtrace --- цепочку вызовов. Остановка меняет время работы: сеть может потерять
соединение, а часть периферии продолжит работать. Учебный опыт проводим на плате без нагрузки.

\subsection{Минимальная сборка ESP-IDF}
Arduino вызывает \code{setup()} и \code{loop()}, а ESP-IDF начинает код с \code{app\_main()}.
Проект имеет корневой CMakeLists, каталог main с регистрацией исходников и sdkconfig.
Хотя Arduino-ESP32 использует ESP-IDF, это разные модели сборки: нельзя просто переопределить
опцию уже собранного загрузчика строкой в скетче.

В активированном терминале ESP-IDF откройте \code{examples/debug\_mean} из репозитория:
\begin{quote}\ttfamily\small
idf.py set-target esp32s3\\
idf.py menuconfig\\
idf.py build\\
idf.py -p COM5 flash monitor
\end{quote}
Замените COM5 своим портом. Фиксируйте версию среды и sdkconfig рядом с проектом.
Файл ELF каждой прошивки нужен для сопоставления машинных адресов и исходника.

\subsection{Ошибка с понятным ожиданием}
\codecap{Arduino-вариант: ожидаем 2,5, получаем 2,0}
\codeinput{l35_debug}

Полный IDF-проект в \code{examples/debug\_mean} содержит ту же функцию в main.c.
Вместо Arduino delay используется \code{vTaskDelay}, вместо Serial --- \code{printf}.
При \code{total=5}, \code{count=2} возвращается 2,0: оба операнда целые, дробь отбрасывается
до преобразования результата в float. Исправьте выражение на \code{(float)total / count};
отдельно определите поведение при нулевом числе элементов.

\subsection{Подключение встроенного USB Serial/JTAG}
GPIO19/20 оставьте свободными, в этой прошивке TinyUSB не используется.
Подключите нативный USB. Используйте OpenOCD Espressif и настройку доступа/USB-драйвер
для своей ОС по официальной инструкции. Обычный системный OpenOCD может не подходить.
Запустите \code{idf.py openocd}; в другом активированном терминале того же проекта --- \code{idf.py gdb}.
В GDB выполните \code{break meanWrong}, \code{continue}, \code{print total}, \code{print count},
\code{next}, \code{bt}. Проверьте, на каком этапе теряется дробная часть.

Для отладочного опыта выберите Debug optimization в menuconfig. Оптимизатор может встроить
функцию или убрать переменную; отсутствие переменной в watch не доказывает, что её нет в исходнике.
После исправления пересоберите, прошейте и отлаживайте именно новый ELF. Не путайте логическую
ошибку этого примера с падением: здесь программа исполняется, но результат неверен.

\begin{trouble}
Точка не ставится: другой ELF или оптимизация. Нет подключения: выбран мост вместо нативного USB,
не установлен драйвер либо прошивка переключила PHY на USB-OTG. Сначала проверьте устройство
и конфигурацию адаптера, затем команды GDB.
\end{trouble}
\begin{summary}
Отладчик показывает причину ошибки, ELF связывает адреса со строками.
Остановка меняет временное поведение; воспроизводимый сценарий важнее случайных наблюдений.
\end{summary}

\begin{selfcheck}
\begin{enumerate}
\item Почему float-результат не спасает целочисленное деление?
\item Зачем нужен ELF этой прошивки?
\item Почему breakpoint мешает сетевому обмену?
\end{enumerate}
\end{selfcheck}
\needspace{20\baselineskip}
\begin{tasks}
\begin{enumerate}
\item \lv{1} Подпишите назначение моста и нативного USB.
\item \lv{1} Воспроизведите 2,00 и запишите ожидаемый результат.
\item \lv{2} Остановитесь в функции и посмотрите аргументы и их типы.
\item \lv{2} Исправьте деление и задайте поведение при count=0.
\item \lv{3} Оформите отчёт: вход, ожидание, факт, стек, причина, правка.
\item \lv{4} Перенесите собственный опыт в ESP-IDF и сохраните версии, sdkconfig, исходники и ELF.
\end{enumerate}
\end{tasks}

\paragraph{Документация для продолжения}\href{https://docs.espressif.com/projects/esp-idf/en/v5.4/esp32s3/api-guides/jtag-debugging/index.html}{Официальное руководство по теме}.
